\section{Bottleneck analysis}
\label{sec:bottlenecks}

\subsection{Ranking on the laptop (best usable configuration)}
For the MMS-300M-CTC, 5-word-cap configuration (median \ttfa\ 5.9\,s), the stage medians are:

\begin{center}\small
\begin{tabularx}{\textwidth}{@{}L{6.2cm}rrY@{}}
\toprule
\textbf{Stage} & \textbf{Median (s)} & \textbf{Share} & \textbf{Nature of the cost} \\
\midrule
ASR after endpoint (offline CTC over the whole turn) & 2.41 & 38\% & compute; avoidable by streaming \\
TTS of the first chunk (under LLM contention) & 2.33 & 37\% & compute + contention \\
LLM to first 5 words (warm cache) & 0.91 & 14\% & compute (prefill + $\sim$8 tokens) \\
Endpoint wait (700\,ms hangover) & 0.70 & 11\% & \emph{pure waiting} \\
\bottomrule
\end{tabularx}
\end{center}
(The shares are of the sum of stage medians, 6.35\,s, which differs slightly from the median \ttfa.)

\textbf{Interpretation.}
\begin{enumerate}
  \item \textbf{ASR finalisation} is the largest term, and almost all of it is avoidable: the audio was
  available while the user spoke. Streaming CTC (E-ASR-4) should reduce it to roughly one chunk
  (0.2--0.4\,s on this CPU\proj).
  \item \textbf{First-chunk TTS} costs 2.3\,s, of which about 0.9\,s is contention. The first chunks are
  about 2.5\,s of audio, which would take about 1.4\,s on an idle CPU (RTF 0.57). While the LLM keeps
  decoding in parallel, the effective RTF rises to 0.8--1.2\meas. Contention control (E-SYS-2) and a
  faster runtime or voice (E-TTS-5/6) attack it.
  \item \textbf{The LLM} is not the latency problem on this hardware once the prompt is cached; it is the
  \emph{quality} problem.
  \item \textbf{The endpoint wait} is small in this configuration only because the other terms are large. On a
  GPU it becomes the dominant term.
\end{enumerate}

\subsection{Projected budgets after optimisation}
The table below is a \emph{projection}\proj\ built from our CPU measurements, the literature numbers in
Section~\ref{sec:components}, and the expected effect of the experiments. It is the hypothesis Week~7
has to confirm or reject, not a result.

\begin{center}\small
\begin{tabularx}{\textwidth}{@{}L{3.2cm}rrY@{}}
\toprule
\textbf{Stage} & \textbf{CPU now}\meas & \textbf{CPU target}\proj & \textbf{How} \\
\midrule
Endpoint wait & 0.70 & 0.30--0.45 & Amharic end-of-turn model + speculative start (E-EP-3/4/5) \\
ASR after endpoint & 2.41 & 0.25--0.40 & Streaming CTC during speech; ONNX/OpenVINO (E-ASR-4/6) \\
LLM to first chunk & 0.91 & 0.6--0.8 & Caching (done), spoken prompt, pinned cores (E-LLM-4/7) \\
TTS first chunk & 2.33 & 0.5--1.0 & Contention control, ONNX VITS or Piper voice (E-SYS-2, E-TTS-5/6) \\
\midrule
\textbf{\ttfa} & \textbf{5.9} & \textbf{1.7--2.7} & ``Acceptable'' tier on a 2016 laptop \\
\bottomrule
\end{tabularx}
\end{center}

\begin{center}\small
\begin{tabularx}{\textwidth}{@{}L{3.2cm}rY@{}}
\toprule
\textbf{Stage} & \textbf{T4 target}\proj & \textbf{Basis} \\
\midrule
Endpoint wait & 0.30--0.45 & Same as CPU (waiting does not get faster on a GPU) \\
ASR after endpoint & 0.05--0.15 & CTC encoders run at RTF $\sim$0.001--0.006 on GPU~\cite{omniasr}; the last streaming chunk only \\
LLM to first chunk & 0.3--0.5 & Gemma~4 E2B/E4B in 4-bit, cached prompt, $\sim$16 tokens for the first 5 words \\
TTS first chunk & 0.2--0.5 & MMS-TTS on GPU, or OmniVoice with 8--16 steps on a short chunk~\cite{omnivoice} \\
Network (laptop $\leftrightarrow$ notebook) & 0.2--0.4 & One response-direction trip plus the tunnel; to be measured (E-SYS-4) \\
\midrule
\textbf{\ttfa} & \textbf{1.1--2.0} & ``Good'' tier, with the endpoint wait and the network now dominant \\
\bottomrule
\end{tabularx}
\end{center}

\begin{finding}[Where Week 7 should spend its time]
On the CPU: ASR streaming and TTS contention first (together $\approx$75\% of \ttfa). On the GPU: the
endpoint wait and the network, because compute stops being the bottleneck. Both paths share the
Amharic end-of-turn work. That is why E-EP-3/5 are P0 even though they save less on the CPU today.
\end{finding}
